\section{Score Matching: Aprendizaje de la función Score}

\subsection{Planteamiento del problema}

En las sesiones anteriores vimos que: (1) la función score describe completamente la distribución de probabilidad; (2) mediante dinámicas de Langevin se puede realizar muestreo utilizando la función score. Entonces, la pregunta clave se convierte en: \textbf{¿cómo aprender la función score a partir de los datos?}

La idea más directa es definir un objetivo de coincidencia de scores:

$$\mathcal{L}_{\text{SM}} = \mathbb{E}_{q(x)}\left[\left\| s_\theta(x) - \nabla_x \log q(x) \right\|^2\right]$$

donde $s_\theta(x)$ es el estimador de score parametrizado por una red neuronal.

\begin{warningbox}{Dificultades del Score Matching directo}
La función de pérdida anterior tiene un problema fundamental: requiere conocer la función score $\nabla_x \log q(x)$ de la distribución real $q(x)$, y es precisamente lo que queremos aprender. Esto parece caer en un ciclo de ``el huevo pone la gallina, la gallina pone el huevo''.
\end{warningbox}

\subsection{Denoising Score Matching}

El insight clave para resolver este problema es que, aunque no conozcamos la función score de la distribución original $q(x)$, si definimos una versión \textbf{con ruido} de la distribución $q_\sigma(x_t)$:

$$q_\sigma(x_t) = \int q(x_t \mid x_0)\, q(x_0)\, dx_0$$

donde $q(x_t \mid x_0) = \mathcal{N}(x_t;\; x_0, \sigma^2 I)$, entonces el \textbf{score condicional} $\nabla_{x_t} \log q(x_t \mid x_0)$ es conocido:

$$\nabla_{x_t} \log q(x_t \mid x_0) = -\frac{x_t - x_0}{\sigma^2}$$

\begin{importantbox}{Pérdida de Denoising Score Matching (DSM)}
Después de una deducción matemática (omitiendo los detalles de las transformaciones de identidad), se puede demostrar que el objetivo de score matching es equivalente a:
$$\mathcal{L}_{\text{DSM}} = \mathbb{E}_{q(x_0)}\, \mathbb{E}_{q(x_t \mid x_0)}\left[\left\| s_\theta(x_t) - \nabla_{x_t} \log q(x_t \mid x_0) \right\|^2\right]$$
Dado que $\nabla_{x_t} \log q(x_t \mid x_0)$ tiene una solución en forma cerrada, ¡esta pérdida es computable!
\end{importantbox}

\begin{knowledgebox}{Proceso de entrenamiento de DSM}
Los pasos de entrenamiento de Denoising Score Matching son extremadamente simples:
\begin{enumerate}
    \item Muestrear un punto de datos $x_0 \sim q(x_0)$ (es decir, tomar una muestra de entrenamiento).
    \item Muestrear ruido $\epsilon \sim \mathcal{N}(0, I)$, obteniendo $x_t = x_0 + \sigma\, \epsilon$.
    \item Calcular el objetivo de score $\nabla_{x_t} \log q(x_t \mid x_0) = -\epsilon / \sigma$.
    \item Minimizar $\|s_\theta(x_t) + \epsilon/\sigma\|^2$.
\end{enumerate}
Se puede observar que esto es esencialmente \textbf{idéntico} al entrenamiento de predicción de ruido en DDPM.
\end{knowledgebox}

\subsection{Equivalencia con el entrenamiento de DDPM}

En DDPM, el objetivo de entrenamiento del predictor de ruido es:

$$\mathcal{L}_{\text{DDPM}} = \mathbb{E}_{x_0, t, \epsilon}\left[\left\| \epsilon_\theta(x_t, t) - \epsilon \right\|^2\right]$$

Mientras que en denoising score matching, el objetivo de entrenamiento de la red score es:

$$\mathcal{L}_{\text{DSM}} = \mathbb{E}_{x_0, t, \epsilon}\left[\left\| s_\theta(x_t, t) + \frac{\epsilon}{\sqrt{1-\bar{\alpha}_t}} \right\|^2\right]$$

La relación entre ambos es $s_\theta(x_t, t) = -\epsilon_\theta(x_t, t) / \sqrt{1-\bar{\alpha}_t}$; es decir, la red score y la red de predicción de ruido difieren únicamente por un factor de escala relacionado con el paso de tiempo.

\begin{importantbox}{Unificación de DDPM y Modelos basados en Score}
El marco de predicción de ruido de DDPM y el marco de score matching basado en score son \textbf{dos caras de la misma moneda}:
\begin{itemize}
    \item Perspectiva de DDPM: entrenar una red neuronal para predecir el ruido añadido a los datos.
    \item Perspectiva basada en score: entrenar una red neuronal para estimar la función score de la distribución con ruido.
    \item Ambas perspectivas son matemáticamente completamente equivalentes, diferenciándose solo por un factor de escala.
\end{itemize}
\end{importantbox}

\subsection{Resumen de la sección}

Score matching resuelve el problema de ``cómo aprender la función score a partir de los datos'' mediante una astuta transformación del objetivo de score marginal incomputable en un objetivo de denoising score matching computable. Este objetivo de entrenamiento es matemáticamente completamente equivalente al entrenamiento de predicción de ruido de DDPM, unificando así dos direcciones de investigación que parecían diferentes.
